9.1 Purpose of the Experiment
The purpose of the BED v3 field experiment is to evaluate whether timing‑optimized behavioral incentives can increase consumer engagement while maintaining reward budget efficiency. As stated in your uploaded document:

“The experiment evaluates the causal impact of BED policies on event frequency, habit formation, and reward spend.”

The experiment is designed to:

measure behavioral lift

quantify cost‑effectiveness

validate CMDP‑generated policies

support causal inference

provide evidence for operational deployment

9.2 Experimental Design Overview
The experiment uses a randomized controlled trial (RCT) with parallel arms:

EXP001 — Control Group  
Default policy (baseline reward logic)

EXP002 — Treatment A (optional)  
Higher responsiveness weighting (habit‑boosting)

EXP003 — Treatment B (optional)  
Lower soft‑budget factor (cost‑saving)

Your uploaded documents emphasize:

“control vs treatment, randomized blocks, measurable outcomes, statistical tests.”

This design satisfies those requirements.

9.3 Randomization Strategy
Users are randomly assigned at onboarding:

𝑃
(
assignment
)
=
1
𝐾
Where K = number of active experiment groups.

Randomization is:

individual‑level

uniform

persistent (stored in Firestore)

immutable (cannot change after assignment)

This ensures internal validity.

9.4 Treatment Conditions
9.4.1 EXP001 — Control
Default policy

Baseline responsiveness

Standard reward multipliers

No habit‑boosting

No budget tightening

9.4.2 EXP002 — Treatment A (optional)
Increased responsiveness weighting

Stronger habit reinforcement

Expected to increase event frequency

9.4.3 EXP003 — Treatment B (optional)
Lower soft‑budget factor

Reduced reward spend

Expected to maintain behavior at lower cost

These treatments align with your CMDP policy parameters.

9.5 Outcome Metrics
Primary Outcome
Event Frequency  
Number of events per user per week.

Secondary Outcomes
Habit Trajectory  
Change in habit state over time.

Return Time  
Days between events.

Reward Spend  
Total points issued.

Cost per Incremental Event (CPIE)

Δ
Reward Spend
Δ
Events
Fairness Metrics  
Parity across segments (e.g., low‑habit vs high‑habit users).

These metrics match the simulation outputs and your uploaded specification.

9.6 Data Collection
Data is collected automatically via:

Firestore event logs

Cloud Function reward logs

Habit state updates

Experiment assignment records

Policy parameters

Each event record includes:

timestamp

user ID

experiment group

reward issued

habit state

recency

policy zone

This ensures complete traceability.

9.7 Statistical Analysis Plan
9.7.1 Estimation Strategy
For each treatment arm:

𝜏
^
=
𝑌
ˉ
𝑡
𝑟
𝑒
𝑎
𝑡
𝑚
𝑒
𝑛
𝑡
−
𝑌
ˉ
𝑐
𝑜
𝑛
𝑡
𝑟
𝑜
𝑙
Where Y is the outcome metric.

9.7.2 Hypothesis Tests
Two‑sample t‑test for event frequency

Mann‑Whitney U for non‑normal distributions

OLS regression with covariates (optional)

Difference‑in‑means for reward spend

Survival analysis for return time

9.7.3 Significance Threshold
𝛼
=
0.05
9.7.4 Power Analysis
Simulation results provide:

expected effect size

variance

required sample size

This satisfies MIT’s requirement for experimental rigor.

9.8 Ethical and Operational Considerations
No personally identifiable information (PII) is used.

All incentives are non‑coercive.

Users may opt out at any time.

Experiment does not disadvantage any group.

Reward budgets are capped.

This aligns with MIT’s IRB expectations.

9.9 Implementation Plan
Step 1 — Assign users to EXP001
Already completed.

Step 2 — (Optional) Activate EXP002/EXP003
Only if testing treatments.

Step 3 — Collect event data
Automatically handled by Cloud Functions.

Step 4 — Analyze results
Using simulation‑aligned metrics.

Step 5 — Integrate findings into BED v3 document
Forms the Results chapter.

9.10 Summary
This experiment plan provides:

a rigorous RCT design

clear treatment conditions

measurable outcomes

statistical analysis methods

ethical safeguards

operational feasibility

It satisfies MIT CTL’s requirements for:

causal inference

experimental validity

reproducibility

academic rigor
